\chapter{模拟、基准与验收}
\section{验收原则}
验证分为数学恒等式、数值求解、接口行为和科学范围四层。伴随内积测试与参考实现比较检查算子；合成真值检查指定模型下的恢复；任务、协议与文件测试检查可调用性；限制清单则说明哪些实验结论不能从测试中推出。

本次自动化结果为 \TestPassed{} 项通过、\TestSkipped{} 项跳过。跳过项为 GPU/CPU 数值一致性测试和官方 MCP SDK 互操作测试，分别因为没有 CUDA/CuPy 和未安装官方 SDK。测试覆盖新旧协议的真实 stdio 子进程、HTTP 请求验证、取消与中断、不可变数据、结果篡改检测、定量拒绝，以及原有数学内核回归。它不是所有客户端、所有操作系统和所有仪器格式的兼容性认证。

\section{谱--时间联合模拟}
构造七个通道、三个组分、600 个时间样点，$\Delta t=0.5$ s。质量通道为 2、15、28、31、44、45、46，响应矩阵见第二章。三个源项包括平滑矩形和高斯型脉冲；每组分使用不同的因果核。每通道另加常数与线性背景，并加入标准差为 0.006 a.u. 的独立 Gaussian 噪声。

使用 \code{channel_pars}，$\lambda_2=0.018$、$\lambda_1=0.0007$、一阶索引差分、非负约束，最多 2500 次迭代，相对驻点阈值 $10^{-5}$。本次在 100 次迭代后满足阈值。三个源项合并 RMSE 为 0.010632，忽略背景的对照为 0.125152；伴随误差约 $5.19\times10^{-15}$。这些数值来自固定随机种子 20261114，不是从实际实验反推的性能指标。

\fig{04_source_A}{.88}{合成组分 A 的恢复。比较真值、联合背景处理和忽略背景的对照；组分标签不对应化学鉴定。}
\fig{04_source_B}{.88}{合成组分 B 的恢复。各组分有独立传输核，联合模型同时利用谱与时间信息。}
\fig{04_source_C}{.88}{合成组分 C 的恢复。背景消元降低了本例中的系统偏差，不能据此断言任意背景都可被唯一分离。}
\fig{05_residual}{.91}{联合拟合后质量 28 通道的带符号残差，单位为 a.u.。这里只展示一个通道的时间曲线，完整残差数组随结果提供；残差较小不等于候选化学结构已确认。}

\section{质量轴与背景模拟}
质量轴测试使用 44 至 46 的均匀网格，共 501 点，三个窄峰与固定 Gaussian 线形。观测相对真值的 RMSE 为 0.044483，恢复 RMSE 为 0.003221，180 次迭代后收敛。它仅验证已知、平移不变线形条件下的一维恢复。

背景测试在 0 至 240 s 区间构造参考通道、压力、目标通道与带正负变化的附加残差。空白窗预设为 $[0,45]$ 和 $[205,240]$ s。45 通道仅用参考通道的残差 RMSE 为 0.027964；加入压力项后为 0.005081。共保留 210 个负残差样点。因为真值按同类模型构造，本例有利于压力项回归；真实实验应额外测试模型失配与独立空白。

\begin{table}[htbp]\centering
\caption{合成实验结果。量纲均为各自模拟信号的 a.u.，不是检出限。}
\begin{tabular}{p{.43\linewidth}rr}\toprule
实验 & 对照 RMSE & 处理后 RMSE\\\midrule
谱--时间源项：忽略背景与联合处理 & 0.125152 & 0.010632\\
质量轴：观测与恢复 & 0.044483 & 0.003221\\
45 背景：参考通道与压力辅助 & 0.027964 & 0.005081\\\bottomrule
\end{tabular}
\end{table}

\section{测量环境与计时约定}
本次 CPU 为 AMD EPYC 9V74，操作系统为 Linux x86\_64；Python 3.13.5、NumPy 2.3.5、SciPy 1.17.0。以环境变量把 BLAS/OpenMP 线程数设为 1，数组使用 float64。暖调用预热两次、测量九次并取中位数；初始化另测五次。固定迭代实验使用 100 次迭代，预热一次、测量三次。

暖调用不含算子分配和核 FFT，不能把它当作文件导入到结果导出的端到端延迟。CPU 调度和缓存状态也使重复计时存在波动；全部样本而不只是比值保存于 \code{benchmarks/ms_results.json}。

\begin{table}[htbp]\centering\small
\caption{相同离散模型下的暖调用中位数。维度按 $(M,J,N,K)$，投影行最后一维为基数 $p$。}
\setlength{\tabcolsep}{4pt}
\begin{tabular}{p{.26\linewidth}p{.25\linewidth}rrr}\toprule
任务 & 维度 & 对照/ms & 优化/ms & 倍数\\\midrule
谱--时间前向 & $(8,3,4096,129)$ & 0.2049 & 0.1124 & 1.82\\
谱--时间伴随 & $(8,3,4096,129)$ & 0.2050 & 0.1342 & 1.53\\
谱--时间前向 & $(24,12,8192,257)$ & 1.7497 & 0.8572 & 2.04\\
谱--时间伴随 & $(24,12,8192,257)$ & 1.6881 & 0.9252 & 1.82\\
谱--时间前向 & $(32,16,32768,513)$ & 12.1626 & 6.3994 & 1.90\\
谱--时间伴随 & $(32,16,32768,513)$ & 11.1867 & 9.3064 & 1.20\\
背景投影 & $(24,4,2048,2)$ & 0.7846 & 0.1408 & 5.57\\
固定 100 次 FISTA & $(12,6,2048,65)$ & 122.7308 & 79.3912 & 1.55\\
\bottomrule\end{tabular}

\end{table}

\fig{12_speedup}{.94}{本次质谱相关实测加速比。每个比值来自记录的中位数；不包含 GPU 数字，也不包含图像任务。}
\fig{13_background_memory}{.85}{背景投影因子的存储量。稠密整体 Q 与分通道因子比较约相差 23.96 倍；该量不是应用程序的实测峰值内存。}

\section{直接法、FFT 与固定迭代}
在 $N=4096,K=31$ 的短核一维案例中，直接法与 FFT 接近，直接法对 FFT 的速度比为 0.984，说明 FFT 并未占优；在 $N=131072,K=1025$ 案例中，相同口径的比值为 6.481。基准同时记录 C++ 与 overlap-add，不能只挑选最快方法后宣称所有输入都会有相同加速。

谱--时间前向的批量 FFT 相对逐组分缓存 FFT 在本次三组维度中约快 1.82、2.04、1.90 倍，伴随约快 1.53、1.82、1.20 倍。分通道投影约快 5.57 倍。固定 100 次 FISTA 的运行约由 122.73 ms 降至 79.39 ms，即 1.55 倍；这两次运行均未以给定极严阈值收敛，因此只能解释为固定迭代预算的速度比较，不是达到相同收敛精度所需时间的比较。

\section{数值一致性与验收边界}
批量与逐组分谱--时间计算在所测样本中输出相同；投影结果最大绝对差约 $4.44\times10^{-16}$。一致性测试证明本次优化没有改变被比较的离散运算，不能代替真实仪器的响应校准。

仍需实验验证的内容包括：单组分与混合物的响应矩阵、核随流量和介质的稳定性、回收率与质量守恒、重复实验精度、背景模型外推误差，以及标定不确定性。本次没有执行这些实验，报告中不提供由模拟直接换算的检出限、准确度或产物选择性。

\section{尚未覆盖的功能}
本版 MS MCP 不是完整的所有类型质谱软件。它不包含供应商二进制读取器、自动 mzML 工作流、多电荷蛋白反卷积、MS/MS 分子鉴定、同位素标记全网络拟合、时变非线性传输，或顺序通道扫描的自动重建。原有通用数学内核可作为进一步实现的基础，但不能把潜在可扩展性写为已经验证的功能。
